\documentclass[11pt]{article}
\usepackage[utf8]{inputenc}
\usepackage{amsmath,amssymb,amsthm}
\usepackage[margin=1in]{geometry}
\usepackage{hyperref}
\usepackage{xcolor}
\usepackage{enumitem}
\usepackage{newunicodechar}
\newunicodechar{ℝ}{\ensuremath{\mathbb{R}}}
\newunicodechar{ℕ}{\ensuremath{\mathbb{N}}}

\title{Tide retrospective:\\\emph{Generic even-monomial 1D template (A1)}}
\author{Daniel Murfet}
\date{18 May 2026}

\begin{document}
\maketitle

\begin{abstract}
Programme-level lift of the 1D quartic / sextic template to arbitrary
$k \ge 1$. For the centred even-monomial potential
$L_k(x) = x^{2k}/(2k)!$, lands the exact closed forms for the
partition function, even moments, and expected values, plus
odd-moment vanishing. After A1, the quartic and sextic seabed
specialisations are mathematically (though not yet syntactically)
five-line corollaries. 310 lines, single session, three trailing
\texttt{ring} / \texttt{field\_simp} cleanups needed. Zero
\texttt{sorry}, zero \texttt{axiom}, zero \texttt{native\_decide}.
Integrability witnesses (\texttt{kth\_integrable\_pow}) are
deliberately deferred to a follow-up tide --- the Gaussian-domination
argument needs case analysis on $k = 1$ vs $k \ge 2$ and is
substantial enough to warrant its own deliberation.
\end{abstract}

\section{Setting}

Today's V2 / Z1 / Z2 trio packaged exact 2D quartic-sextic
asymptotics; A1 lifts the underlying 1D template to arbitrary $k$.
N1 (the generic-$k$ J-function tide) had already established the
\emph{partition function in raw-integral form}; A1 fills the rest of
the template (moments, odd-vanishing, partition in
\texttt{partitionFunction}-form, expected values), matching the shape
of the quartic and sextic files structurally.

The strategic value: after A1 lands, downstream programmes using
pure-monomial potentials at arbitrary $k$ (or 2D separable
combinations of such potentials) can build on a parametric template
rather than ad-hoc quartic / sextic specialisations. The today's
quartic-sextic 2D story has effectively prototyped what the template
enables.

\section{The thing formalised}

For $k \ge 1$, $j \in \mathbb N$, and $t > 0$:

\begin{itemize}[leftmargin=1.5em]
\item \emph{Half-line moment}: $\int_0^\infty x^{2j}\,e^{-tx^{2k}/(2k)!}\,dx
  = \tfrac{1}{2k}\cdot ((2k)!/t)^{(2j+1)/(2k)}\cdot \Gamma((2j+1)/(2k))$.
\item \emph{Full-line moment} (\texttt{kth\_moment\_even}):
  $\int_{\mathbb R} x^{2j}\,e^{-tx^{2k}/(2k)!}\,dx
  = \tfrac{1}{k}\cdot ((2k)!/t)^{(2j+1)/(2k)}\cdot \Gamma((2j+1)/(2k))$.
\item \emph{Odd-vanishing} (\texttt{kth\_moment\_odd}):
  $\int x^{2j+1}\,e^{-tx^{2k}/(2k)!}\,dx = 0$.
\item \emph{Partition in \texttt{partitionFunction}-form}:
  $Z_{L_k}(t) = \tfrac{1}{k}\cdot ((2k)!/t)^{1/(2k)}\cdot \Gamma(1/(2k))$.
\item \emph{Expected values} (even, odd): moments divided by the
  partition; even case is $((2k)!/t)^{j/k}\cdot
  \Gamma_{\rm num}/\Gamma_{\rm den}$, odd case vanishes.
\end{itemize}

Named definition \texttt{kthPotential} + \texttt{kthPotential\_apply}
\texttt{@[simp]}. File: sibling to the quartic and sextic files in
\texttt{Laplace/OneD/}.\footnote{Specifically
\texttt{Laplace/OneD/MonomialPotential.lean}; per GPT's preference
for descriptive over parametric naming.}

\section{Proof strategy}

The half-line moment is a direct call to Mathlib's
\texttt{integral\_rpow\_mul\_exp\_neg\_mul\_rpow} at $p = 2k$,
$q = 2j$, $b = t/(2k)!$, with the standard cast-handling massage
to convert Nat-power form to rpow form for the application. The
full-line moment follows by \texttt{integral\_comp\_abs} on the
even-symmetric integrand; the factor of 2 absorbs into $2/(2k) =
1/k$. Odd-vanishing is via odd reflection symmetry against the
even-symmetric measure, identical template to
\texttt{quartic\_moment\_odd}. Partition is the $j = 0$
specialisation. Expected values are the moment / partition ratio,
with an \texttt{rpow\_add} step to combine
$((2k)!/t)^{(2j+1)/(2k) - 1/(2k)} = ((2k)!/t)^{j/k}$.

\section{Roadblocks and resolutions}

Three trailing tactic mismatches at the bottom of the longer proofs:

\begin{enumerate}[leftmargin=1.5em]
\item In the half-line integral proof,\footnote{Theorem
  \texttt{integral\_pow\_mul\_exp\_neg\_kth\_Ioi}.} an \texttt{hg}
  sub-lemma used \texttt{push\_cast; ring}, but \texttt{push\_cast}
  alone closed the goal so \texttt{ring} errored with ``No goals to
  be solved''. The asymmetric \texttt{hgneg} needed \texttt{ring}
  because \texttt{push\_cast} stopped one step short. Fix:
  \texttt{push\_cast} alone for the positive case, \texttt{ring}
  after \texttt{push\_cast} for the negative.
\item Same pattern in two further proofs:\footnote{The full-line
  \texttt{kth\_moment\_even} and the expected-value lemma both ended
  with \texttt{field\_simp; ring}; in both, \texttt{field\_simp}
  closed the goal alone.} \texttt{field\_simp} closed the goal
  alone in each. Fix: drop the trailing \texttt{ring}.
\item In the \texttt{partitionFunction\_kthPotential} bridge, an
  \texttt{hexp} helper for the exponent reduction used the form
  \texttt{(2 * 0 + 1 : ℝ) / ...} but the post-\texttt{kth\_moment\_even}
  goal had Nat-cast \texttt{↑0} still visible. The \texttt{rw}
  failed to match. Fix: replace the explicit \texttt{hexp} +
  \texttt{rw} with a single \texttt{norm\_num}, which normalises
  \texttt{↑0} and simplifies \texttt{2 * 0 + 1 = 1} in one pass.
\end{enumerate}

All four were local mid-proof cleanups, not architectural problems.
No GPT mid-proof consult was needed; the deliberation at Step 2 set
the scope and naming correctly.

\section{What was Mathlib, what was new}

\paragraph{From Mathlib.} The master lemma
\texttt{integral\_rpow\_mul\_exp\_neg\_mul\_rpow},
\texttt{integral\_comp\_abs}, \texttt{integral\_neg\_eq\_self}, the
\texttt{Real.rpow} algebra suite, \texttt{Real.Gamma\_pos\_of\_pos},
and standard \texttt{ring} / \texttt{field\_simp} /
\texttt{push\_cast} / \texttt{norm\_num}.

\paragraph{From the seabed.} Only N1's \texttt{kth\_partition} for
the raw integral identity is in the spirit; A1 inherits the same
proof template (via Mathlib's master lemma at $q = 0$) and extends
to general $q = 2j$. The quartic and sextic seabed lemmas are not
load-bearing for A1 (A1 is generic enough to subsume both, even if
the actual refactor of the quartic / sextic files is deferred to a
follow-up tide).

\paragraph{New here.} The named \texttt{kthPotential} definition,
the full eight-theorem template at parametric $k$, and the bridge
between raw-integral and \texttt{partitionFunction} /
\texttt{gibbsExpectation} forms.

\paragraph{Deferred.} Integrability witnesses
(\texttt{kth\_integrable\_pow}, \texttt{kth\_integrable\_pow\_pot}).
The quartic and sextic versions use a Gaussian-domination argument
via square completion of $tx^{2k}/(2k)! - x^2$ for $k = 2, 3$
specifically; a generic-$k$ version needs a uniform constant
$C(k, t)$ that splits cleanly into the $k = 1$ (direct Gaussian)
and $k \ge 2$ (square-completion with parametric constant) cases.
That's substantial enough to warrant its own tide.

\section{Lessons}

\paragraph{Proceeding without GPT mid-proof works for mirror-shaped
extensions.} GPT's pre-proof consult at Step~2 was load-bearing: it
voted Candidate A, named the naming triplet, and explicitly
recommended deferring the quartic / sextic refactor to a separate
tide. Beyond that, A1's proof is mechanical mirror of the quartic /
sextic templates with rpow/cast plumbing parametric in $k$. The
deliberation-once + execute-deterministically pattern scales.

\paragraph{\texttt{field\_simp} and \texttt{push\_cast} are
``finishing'' tactics in this regime.} Four of the proof bodies in
this file ended with what was meant to be \texttt{push\_cast; ring}
or \texttt{field\_simp; ring}, but the closing \texttt{ring} was
redundant in three of those four cases --- the prior tactic closed
the goal. Worth recording: after explicit \texttt{rpow\_add /
rpow\_neg} normalisations, \texttt{ring} is often unnecessary. If
\texttt{ring} errors with ``No goals to be solved'', drop it.

\paragraph{Defer integrability when it forks the proof.} The generic
Gaussian-domination argument forces a case split on $k = 1$ vs
$k \ge 2$ that the quartic / sextic specialisations don't need. By
deferring the integrability layer, A1 stays as a clean uniform
parametric lift; the integrability tide can then specialise as
needed without contaminating the headline theorems.

\section{Follow-ups}

\begin{itemize}[leftmargin=1.5em]
\item \emph{Integrability tide.} Land
  \texttt{kth\_integrable\_pow} and \texttt{kth\_integrable\_pow\_pot}
  via the Gaussian-domination argument, case-split on $k = 1$ vs
  $k \ge 2$. ${\sim}80$ lines.
\item \emph{Refactor quartic and sextic files as specialisations.}
  After A1 and the integrability follow-up land, rewrite the quartic
  and sextic files so that their headline theorems become $k = 2$
  and $k = 3$ specialisations of the generic ones. Net deletion
  ${\sim}300$ lines. Per GPT's guidance, this is a separate tide.
\item \emph{2D analogues at arbitrary $k_1, k_2$.} The 2D
  quartic-sextic story today proved closed for $(k_1, k_2) = (2, 3)$;
  with A1 in hand, the same V2/Z1/Z2-style packaging extends to
  arbitrary $(k_1, k_2)$ separable monomial potentials.
\end{itemize}

\section*{Acknowledgements}

GPT-5.5~Pro's Step-2 consult contributed the vote for Candidate A
(full template over skinnier variants), the naming triplet, the
\texttt{MonomialPotential.lean} file name, and the explicit
``defer the quartic / sextic refactor'' guidance. Full response
preserved alongside the tide log in the primer project's
\texttt{gpt\_responses\_a1/} folder.

\end{document}
